\textbf{Construcción de requisitos (cláusula 5.2.4).} Un requisito bien formado expresa una necesidad del sistema de interés y no una simple preferencia de interfaz o una descripción de pasos de implementación. La construcción debe vincular el requisito con un problema o intención de un interesado, expresar una capacidad o condición relevante y permitir comprobar objetivamente su cumplimiento. La diferencia entre decir «Spotify debe ser mejor» y especificar «la búsqueda devolverá resultados para un título del catálogo» está en el resultado observable. Además, la redacción debe distinguir restricciones impuestas de decisiones de diseño. El hecho de que una característica exista hoy no demuestra por sí mismo que sea una necesidad; la decisión de conservarla requiere justificación \cite{iso29148,pohl2025}.

\textbf{Calidad de un requisito individual (cláusula 5.2.5).} La norma identifica nueve características. \emph{Necesario}: su ausencia dejaría una necesidad sin resolver. \emph{Apropiado}: el detalle corresponde al nivel de sistema o software sin fijar soluciones innecesarias. \emph{No ambiguo}: conduce a una interpretación coherente. \emph{Completo}: contiene la información indispensable para comprender la obligación. \emph{Singular}: se concentra en una función o cualidad, aunque pueda tener condiciones. \emph{Factible}: su realización es viable con recursos y límites razonables. \emph{Verificable}: admite demostración, inspección, análisis o prueba. \emph{Correcto}: conserva la intención de la necesidad original. \emph{Conforme}: sigue el formato convenido. Estas propiedades se relacionan: agregar una cifra no convierte automáticamente una propuesta en correcta o necesaria. Por ejemplo, un límite de tres segundos puede facilitar una prueba, pero antes debería justificarse con necesidades reales y condiciones de uso \cite{iso29148}.

\textbf{Calidad del conjunto (cláusula 5.2.6).} Revisar cada oración por separado no basta. El conjunto debe ser \emph{completo}, porque recoge las necesidades dentro del alcance declarado sin asuntos sin resolver; \emph{consistente}, porque no contiene mandatos contradictorios o duplicados y mantiene vocabulario y unidades compatibles; \emph{factible}, por poder realizarse de forma conjunta; \emph{comprensible}, de modo que los involucrados entiendan el sistema esperado; y \emph{validable}, porque satisfacer el conjunto debería responder a las necesidades del interesado. En Spotify, un requisito que permite descargar canciones a toda cuenta entraría en conflicto con otro que condiciona esa función a un plan Premium; la revisión del conjunto detecta el problema que una evaluación aislada podría omitir \cite{iso29148}.

\textbf{Lenguaje controlado (cláusula 5.2.7).} Conviene evitar superlativos como «el mejor»; apreciaciones subjetivas como «intuitivo»; adjetivos sin escala como «rápido»; pronombres sin referente; conectores que admiten lecturas alternativas; frases abiertas («entre otros», «etc.»); comparaciones sin base; cláusulas de escape («si es posible»); absolutos («siempre», «nunca», «todos») y referencias incompletas. La norma recomienda expresar \emph{qué} necesita el sistema, no imponer \emph{cómo} se construye. La selección cuidadosa de palabras reduce interpretaciones incompatibles; la métrica, la condición y el criterio de aceptación transforman una expectativa en una obligación examinable \cite{iso29148}.

\textbf{Diferencia entre necesidad y solución.} Una especificación técnica no debe confundirse con una descripción de la interfaz en su versión actual. Por ejemplo, «la aplicación utilizará un botón verde ubicado a la izquierda» fija innecesariamente una alternativa visual, mientras «el usuario podrá identificar la función para crear una lista» expresa una capacidad. El segundo enunciado deja margen a distintas soluciones sin perder su propósito. La excepción aparece cuando un interesado impone expresamente una restricción de diseño: en ese caso el origen y la justificación deben quedar documentados. En un servicio comercial como Spotify, también es necesario reconocer que hay reglas de suscripción que condicionan qué capacidades quedan habilitadas. El requisito debe conservar explícitamente esa condición para no prometer funciones fuera de alcance \cite{iso29148,spotify-playlists}.

\textbf{Distinción entre verificabilidad y validación.} En esta actividad hay dos preguntas distintas. La revisión de la frase identifica si el requisito está bien escrito, por ejemplo si un evaluador puede medir una cantidad con una unidad definida. La evaluación de necesidades examina si esa exigencia tiene sentido para quienes utilizarán el sistema. Un umbral de 120 segundos para aprender a crear una lista puede ser perfectamente medible y, al mismo tiempo, no representar las expectativas de usuarios reales. La elección del tamaño de muestra, los contextos de uso y la tolerancia al error necesita sustento; no nace automáticamente de la norma. Por ello los objetivos cuantitativos de este informe son hipótesis de aceptación que tendrían que discutirse con usuarios y corroborarse con ensayos, no propiedades demostradas de Spotify \cite{iso29148,iso25010}.

\textbf{Precisión del alcance y trazabilidad.} La variedad de planes, dispositivos y condiciones de red afecta la lectura de los requisitos. Una operación posible en escritorio podría tener límites diferentes en un teléfono, mientras la descarga musical está relacionada con funciones Premium. Sin una frontera explícita, dos personas podrían evaluar el mismo enunciado en entornos incompatibles y obtener resultados contradictorios sin que ninguna haya hecho mal la prueba. La trazabilidad permite unir una necesidad con el identificador del requisito, su fuente, el criterio de aceptación y la evidencia de ejecución. Para R-01, por ejemplo, la captura de una lista recién creada da respaldo visual al resultado, pero no demuestra que el proceso sea fiable después de muchas repeticiones. Una captura, una medición y una entrevista cumplen funciones diferentes y no deben sustituirse indistintamente \cite{iso29148,spotify-playlists,spotify-premium}.

\textbf{Aplicación crítica a Spotify.} Se separan funciones sustentadas en la ayuda oficial de metas de calidad todavía no medidas. La biblioteca, la edición de playlists y la sesión privada son funciones documentadas, mientras que la rapidez, estabilidad, facilidad de aprendizaje y sincronización han sido formuladas como objetivos de evaluación. Por ello no se declara que Spotify ya satisfaga los umbrales propuestos. La norma ISO/IEC 25010:2023 facilita clasificar cada meta de calidad en una característica y una subcaracterística, pero no fija por sí sola la cifra de aceptación para este caso: esa cifra se negocia y se valida con usuarios y condiciones de prueba \cite{iso25010,pohl2025}.


\textbf{De la métrica a la prueba.} Una métrica expresa qué propiedad será observada y un umbral indica desde qué resultado se considera aceptable. El requisito de búsqueda, por ejemplo, contiene tiempo de respuesta, número de ensayos y condición de conectividad; aun así, un protocolo riguroso debe fijar dispositivo, versión del cliente, forma de medir y datos de consulta. En una prueba de aprendizaje, además del tiempo hay que definir qué significa finalizar correctamente la tarea y evitar contar resultados de participantes inexistentes. Para privacidad no basta mostrar el interruptor de sesión privada: hay que comprobar el efecto desde otra cuenta autorizada y respetar las condiciones de consentimiento. Por ello el documento separa criterios de evaluación de resultados, conserva pendientes explícitos y no convierte inferencias en evidencias. Esta separación es una forma de aplicar la verificabilidad y la corrección a la vez \cite{iso29148,iso25010}.
